%% !TEX root = manuscript.tex

\chapter{Mission 1 : la reprise des données vers SPO}
\minitoc
\newpage

Ce chapitre présente le chantier le plus lourd de mon année : reconstituer, dans le nouvel ERP,
le socle de données du groupe. Il s'est déroulé en deux temps. Le premier semestre a porté sur
les structures de gestion, c'est-à-dire les opérations, leurs tranches et leurs budgets,
migrées depuis la base historique Grimmo. Le second a porté sur les tiers, c'est-à-dire les
clients acquéreurs, repris depuis le CRM Vadimm.

Ces deux missions partagent la même méthode, mais elles ne posent pas la même question.
Injecter des structures suppose de comprendre un modèle de données inconnu. Injecter des
personnes suppose de décider ce qui fait que deux lignes désignent la même personne, et cette
seconde question s'est révélée beaucoup plus délicate que la première.

\section{La migration des opérations}

Cette mission, conduite au premier semestre, est exposée en détail dans mon mémoire de
mi-parcours. Je n'en reprends ici que ce dont la suite a besoin : la méthode, les deux
contraintes que l'ERP impose, et l'état dans lequel le socle se trouvait à la fin du semestre.

\subsection{Méthode et industrialisation des flux}

Face aux 283 opérations du périmètre, une saisie manuelle intégrale était exclue. Nous avons
donc retenu une stratégie hybride : une injection automatisée pour les gros volumes
structurels, via les interfaces de l'éditeur, et un traitement manuel pour les cas historiques
qui ne pouvaient pas entrer dans les gabarits. Vouloir tout automatiser aurait conduit à écrire
des règles pour des cas uniques, ce qui coûte finalement plus cher que de les traiter à la
main.

J'ai d'abord analysé la documentation technique de l'API de Salvia~[2], puis testé
manuellement les \emph{endpoints} pour comprendre la mécanique de création d'une donnée. Nous
avons ensuite déployé le module DirectAPI, dont j'ai construit le pipeline d'alimentation sous
la forme de maquettes Excel en trois étapes : \textbf{extraction} de la base source par MyReport
Builder, \textbf{transformation et mapping} par formules RechercheX, puis \textbf{chargement}
via DirectAPI. C'est la transformation qui concentre l'essentiel du travail, car la nomenclature
de SPO est stricte : un libellé de poste budgétaire ou un type de tranche non reconnu fait
rejeter la ligne entière.

J'ai sécurisé le processus par une validation itérative : 3 opérations pilotes, puis un lot de
12, avant de lancer la migration de masse. Cette progression par paliers rend les erreurs
lisibles, puisqu'un rejet porte sur un petit nombre de lignes que l'on peut examiner une à une.
C'est cette méthode que j'ai reconduite au second semestre.

\subsection{Deux contraintes de l'ERP, et la documentation qui en est sortie}

La première contrainte est l'ordre strict d'injection, une règle d'intégrité référentielle que
la documentation de l'éditeur n'explicitait pas. Je l'ai découverte en analysant les codes
d'erreur renvoyés par l'API : il faut créer l'opération, puis les tranches de travaux, puis les
tranches commerciales, puis le budget. La figure~\ref{fig:sequence} représente cette séquence.
Toute injection qui ne la respecte pas échoue, le système ne trouvant pas l'objet parent.

\begin{figure}[!htbp]
\centering
\includegraphics[width=0.65\textwidth]{fig_2_1_sequence_injection.png}
\caption{La séquence d'injection imposée par l'intégrité référentielle de l'ERP}
\label{fig:sequence}
\end{figure}

J'ai alors rédigé une documentation interne cartographiant cette séquence et les prérequis de
chaque objet, tout simplement parce que ces règles n'existaient nulle part ailleurs que dans
mes notes de test.

La seconde contrainte tient aux identifiants. Invisibles dans l'interface, ils changent à
chaque mise à jour : l'API ne remplace pas un budget, elle clôture la version précédente et en
crée une nouvelle. Ne pouvant me fier à aucun identifiant fixe, j'ai conçu un requêtage
dynamique en trois temps : un appel GET pour lister les versions existantes, un tri pour retenir
la plus récente, puis un appel POST sur cet identifiant.

\subsection{Le déploiement en « coquilles vides » et le bilan du semestre}

Les utilisateurs devaient pouvoir facturer dès le basculement, notamment pour le lien avec la
Gestion Électronique des Documents (GED). Nous avons donc adopté une stratégie de déploiement
en « coquilles vides », qui dissocie le contenant, c'est-à-dire la structure de l'opération, du
contenu, c'est-à-dire les données budgétaires détaillées. Concrètement, nous avons injecté
d'abord le code opération et le libellé, en reportant les budgets à plus tard.

\begin{longtable}[]{@{}
  >{\raggedright\arraybackslash}p{(\columnwidth - 2\tabcolsep) * \real{0.4000}}
  >{\raggedright\arraybackslash}p{(\columnwidth - 2\tabcolsep) * \real{0.6000}}@{}}
\toprule\noalign{}
\begin{minipage}[b]{\linewidth}\raggedright
Objet du modèle SPO
\end{minipage} & \begin{minipage}[b]{\linewidth}\raggedright
État à la fin du premier semestre
\end{minipage} \\
\midrule\noalign{}
\endhead
\bottomrule\noalign{}
\endlastfoot
Opérations & Créées pour l'ensemble du périmètre, 283 dossiers \\
Tranches de travaux & Créées \\
Tranches commerciales & Création en cours \\
Budgets détaillés & Reportés, stratégie des « coquilles vides » \\
Historique des factures & Non repris \\
\end{longtable}

Sur un projet de reprise, l'avancement se mesure à la part du modèle cible réellement couverte.
Créer 283 opérations sans leurs budgets nous a permis de bien avancer sur la migration et,
surtout, a permis aux autres collaborateurs de commencer à utiliser le logiciel pour leurs
opérations.

\section{La reprise des tiers}

Conduite au second semestre, cette mission prolonge directement la précédente. Elle constitue
le premier pas vers la communication entre l'ERP et le CRM, puisqu'elle transfère le
référentiel client, la synchronisation permanente sans double saisie restant à construire. Elle
change toutefois de nature : il ne s'agit plus d'injecter des structures de gestion, mais des
personnes, avec ce que cela implique de risques de doublons et d'exigence de qualité.

\subsection{Contexte et enjeux}

\subsubsection{Du CRM Vadimm à l'ERP SPO : le périmètre des données clients}

Après la création des opérations et de leurs tranches, il restait à reprendre les données
commerciales, et en premier lieu les clients. Dans SPO, ceux-ci sont désignés par le terme
générique de \emph{tiers} : la même entité de base représente aussi bien un acquéreur
particulier qu'une société, et se voit attribuer un ou plusieurs rôles. Reprendre les tiers,
c'était donc alimenter le socle sur lequel toute la gestion commerciale allait s'appuyer.

Le périmètre consistait à importer l'ensemble des clients ayant déjà acheté chez Angelotti. La
matière première venait du CRM Vadimm, sous la forme de deux extractions complémentaires :
celle des ventes enregistrées, qui porte 1~557 lignes, et celle des coordonnées des acquéreurs
et co-acquéreurs, qui en porte 1~602.

La difficulté ne tenait pas au volume mais à la nature de la donnée. Une ligne de l'extraction
ne décrit pas un client, elle décrit une vente : un acquéreur qui a acheté trois lots y
apparaît trois fois, et rien dans le fichier ne le signale. Transformer un fichier de ventes en
référentiel de personnes suppose donc de reconstituer l'identité de chacune, ce que le système
source ne faisait pas à ma place.

\subsubsection{L'enjeu central : créer sans dupliquer}

Tout l'enjeu se résume à une contrainte : créer chaque client une fois et une seule. Or deux
sources de duplication se cumulaient.

La première est interne au fichier. Chaque ligne correspondant à une vente, un client fidèle y
figure autant de fois qu'il a acheté. Importer le fichier tel quel aurait créé autant de fiches
que d'achats, ce qui aurait rendu impossible la reconstitution de son historique.

La seconde est externe. Au moment de préparer l'import, SPO n'était pas vide : environ 13~000
tiers y avaient déjà été créés par saisie ou par d'autres canaux. Injecter les 1~602 lignes
sans les confronter à cet existant aurait produit des doublons, opposant une fiche importée et
une fiche saisie, toutes deux légitimes.

Une reprise manuelle était donc exclue. Vérifier ligne à ligne qu'un acquéreur n'existe ni
ailleurs dans le fichier ni déjà dans un référentiel de 13~000 tiers représentait un volume de
contrôles incompatible avec les délais, et un risque d'erreur trop élevé. J'ai donc construit un
fichier de travail qui automatise ces contrôles par formules, de sorte que la décision
d'importer une ligne soit calculée plutôt que prise au cas par cas.

\subsection{Conception du fichier de travail}

\subsubsection{Une architecture en trois temps}

Le fichier ne se présente pas comme un tableur de saisie mais comme une chaîne de traitement,
dans la lignée des maquettes d'injection du premier semestre : les données entrent d'un côté
sous leur forme brute et ressortent de l'autre sous la forme attendue par SPO, sans qu'aucune
donnée soit recopiée à la main. La figure~\ref{fig:chaine} en présente les trois étages.

\begin{enumerate}
\def\labelenumi{\arabic{enumi}.}

\item
  \textbf{Les feuilles sources}, en lecture seule : \texttt{VADIMM} porte les 1~557 lignes de
  ventes, \texttt{Builder\ coordonnée} les 1~602 lignes de coordonnées, et \texttt{SPO}
  l'extraction des tiers déjà existants, qui sert de référentiel de comparaison.
\item
  \textbf{La feuille pivot \texttt{BuilderTiers}}, cœur du dispositif : elle reçoit
  l'intégralité des lignes sources et y adjoint les colonnes de calcul qui déterminent, pour
  chacune, s'il s'agit d'un tiers nouveau, d'un doublon interne ou d'un tiers déjà présent dans
  SPO.
\item
  \textbf{Les six feuilles de sortie}, une par fichier d'import attendu par SPO, entièrement
  générées par formules et destinées à être exportées en CSV.
\end{enumerate}

\begin{figure}[!htbp]
\centering
\includegraphics[width=0.65\textwidth]{fig_2_2_chaine_traitement.png}
\caption{La chaîne de traitement à trois étages du fichier de travail}
\label{fig:chaine}
\end{figure}

\subsubsection{Le modèle de données des tiers dans SPO}

Avant d'automatiser quoi que ce soit, il fallait comprendre comment SPO représente un tiers. Je
l'ai fait par lecture de la documentation, puis par tests unitaires de création. Deux
caractéristiques ont structuré tout le reste du travail.

La première est que SPO distingue trois \textbf{natures} de tiers, et que la nature commande la
manière dont l'identité est portée. Une \textbf{personne morale} est une société, identifiée par
sa raison sociale et son SIRET. Une \textbf{personne physique} est un individu, identifié par
son nom, son prénom et sa date de naissance. Un \textbf{couple}, enfin, est une entité
particulière : le tiers n'est ni l'un ni l'autre des conjoints, il les regroupe, et les deux
personnes qui le composent sont enregistrées comme ses \emph{représentants}. Cette dernière
nature est de loin la plus représentée dans nos données, l'essentiel de nos ventes se faisant à
des couples acquéreurs.

La seconde est qu'un tiers n'est pas un enregistrement unique mais un ensemble d'objets liés.
La fiche principale ne porte que l'identité, tandis que l'adresse, le courriel, le téléphone,
les représentants et les rôles sont autant d'objets rattachés, avec chacun son fichier
d'import. C'est ce qui explique que mon fichier produise six feuilles de sortie et non une
seule. Cela explique aussi une asymétrie qui s'est révélée décisive : adresses, courriels et
téléphones se rattachent directement à une personne physique ou morale, mais pour un couple ils
se rattachent aux représentants. Cette différence est devenue la principale contrainte
technique de l'import, comme je l'expose plus loin.

\subsection{Les trois mécanismes automatisés}

La feuille pivot repose sur trois mécanismes qui s'enchaînent : identifier les doublons
internes, écarter les tiers déjà créés dans SPO, puis distribuer les lignes retenues vers les
six feuilles de sortie. Ce sont eux qui distinguent un fichier de mise en forme d'un véritable
outil de reprise, puisqu'ils déplacent la décision d'importer une ligne de l'opérateur vers une
règle écrite.

\subsubsection{La déduplication par clés d'unicité différenciées}

Dédupliquer suppose de répondre à une question qui n'est pas technique : à partir de quoi
décide-t-on que deux lignes désignent la même personne ? C'est le point sur lequel j'ai le plus
hésité, et celui qui commande la qualité du résultat.

Une clé unique appliquée à tout le fichier ne pouvait pas convenir, précisément parce que les
trois natures de tiers ne s'identifient pas de la même façon. Deux sociétés peuvent porter des
raisons sociales très proches sans être la même entité, alors qu'il est très improbable que
deux lignes partageant nom, prénom, date de naissance et adresse désignent deux personnes
distinctes. J'ai donc retenu \textbf{trois clés d'unicité distinctes}, une par nature, calculées
chacune dans sa propre colonne :

\begin{enumerate}
\def\labelenumi{\arabic{enumi}.}

\item
  \textbf{Pour un couple} : le nom et le prénom de l'acquéreur principal, désigné ADV dans nos
  extractions, complétés du nom et du prénom du co-acquéreur, désigné CO. C'est la combinaison
  des deux identités qui fait l'unicité, et non chacune prise isolément.
\item
  \textbf{Pour une personne physique} : le nom, le prénom, la date de naissance et l'adresse. La
  date de naissance sépare les homonymes, l'adresse ajoute un second point de comparaison.
\item
  \textbf{Pour une personne morale} : la raison sociale et le SIRET, identifiant légal de
  l'établissement, la raison sociale prenant le relais lorsque le SIRET fait défaut.
\end{enumerate}

Chaque colonne calcule ensuite un \textbf{rang} : la première occurrence d'un tiers reçoit le
rang 1, les suivantes les rangs 2, 3, et ainsi de suite. Seules les lignes de rang 1 sont
retenues, les autres correspondant à des achats successifs d'un client déjà compté. Ce
mécanisme de rang présente un avantage sur une suppression pure et simple : aucune ligne n'est
effacée, les doublons restent visibles, et l'on peut à tout moment vérifier pourquoi une ligne
n'a pas été importée.

Ce choix de trois clés différenciées est le point méthodologique le plus important de la
mission. Chercher une clé universelle aurait conduit soit à fusionner des tiers distincts, soit
à en dupliquer, et dans les deux cas l'erreur n'aurait été découverte que bien plus tard, dans
l'usage quotidien de l'ERP.

\subsubsection{L'exclusion des tiers déjà présents dans SPO}

Pour chaque ligne de la feuille pivot, trois formules de recherche interrogent l'extraction des
tiers déjà créés dans SPO : les \textbf{personnes morales} par leur raison sociale, les
\textbf{personnes physiques} par leur nom et leur prénom, et les \textbf{couples} par les quatre
noms qui les composent, \textbf{en testant les deux ordres possibles} entre l'acquéreur et le
co-acquéreur.

Lorsqu'une de ces formules renvoie un code tiers de SPO, du type \texttt{T00012666}, la ligne
est exclue de l'import. \textbf{40 tiers} ont ainsi été détectés comme déjà existants, soit une
proportion faible rapportée aux 1~474 tiers uniques du fichier, mais quarante doublons de moins
dans l'ERP.

La figure~\ref{fig:entonnoir} récapitule l'effet cumulé de ces deux mécanismes. Des 1~602
lignes de départ, la déduplication interne en écarte 128 qui correspondent à des achats
répétés, ce qui ramène le total à 1~474 tiers uniques ; l'exclusion en retire 40 de plus, pour
\textbf{1~434 tiers à créer}. Chaque étage correspond à une règle explicite, écrite dans une
colonne, donc vérifiable et reproductible.

\begin{figure}[!htbp]
\centering
\includegraphics[width=0.65\textwidth]{fig_2_3_entonnoir_volumetrie.png}
\caption{De 1~602 lignes de départ à 1~434 tiers créés : l'effet des deux filtres successifs}
\label{fig:entonnoir}
\end{figure}

\subsubsection{La propagation vers les feuilles de sortie}

Les six feuilles de sortie n'ont ni la même population ni le même nombre de lignes. Celle des
tiers reçoit un enregistrement par tiers créé ; celle des adresses ne concerne que les
personnes physiques et morales ; celle des courriels seulement celles dont le courriel est
renseigné ; celle des représentants deux lignes par couple, plus une par société. Il fallait
donc donner à chaque feuille sa propre numérotation.

Le mécanisme retenu attribue à chaque ligne de la feuille pivot un \textbf{numéro séquentiel par
destination} : une ligne qui n'a pas vocation à alimenter une feuille ne reçoit pas de numéro
pour elle, les autres sont numérotées de 1 à \emph{n}. Chaque feuille de sortie n'a plus qu'à
chercher, pour sa ligne \emph{i}, celle de \texttt{BuilderTiers} portant le numéro \emph{i} dans
la colonne qui la concerne.

Ce détour rend les six feuilles indépendantes tout en les gardant synchronisées avec une source
unique. Si une ligne est corrigée ou si une exclusion apparaît, les six numérotations se
recalculent et les fichiers restent cohérents entre eux.

\subsection{De la préparation à l'injection}

\subsubsection{Le contournement par API : 872 requêtes PATCH pour les couples}

C'est ici qu'est apparue la contrainte la plus structurante du projet. Le format d'import CSV
ne permet pas de porter les courriels et les téléphones des couples : pour un couple, ces
coordonnées ne sont pas des attributs du tiers mais de ses représentants, et les fichiers
d'import de coordonnées ne savent viser qu'un tiers. Sur les 881 couples créés, les CSV
pouvaient donc produire l'identité et les représentants, mais pas les coordonnées de contact.
Autrement dit, plus de six tiers sur dix seraient arrivés dans l'ERP sans adresse électronique
ni téléphone.

Saisir ces coordonnées à la main représentait plusieurs centaines de saisies et reconduisait
exactement le risque d'erreur que le fichier automatisé visait à supprimer. J'ai donc cherché
une autre voie, et je me suis tournée vers l'API de SPO, déjà pratiquée au premier semestre.
Elle expose une opération de modification partielle au format \textbf{JSON Patch}~[3], qui
décrit une modification non pas en renvoyant l'objet entier, mais sous la forme d'opérations
élémentaires indiquant l'action, le chemin de la propriété visée et la valeur à appliquer.
C'était exactement ce qu'il me fallait : ajouter un courriel et un téléphone sur un représentant
déjà créé, sans toucher au reste de sa fiche.

J'ai construit le corps de chaque requête à partir de la feuille \texttt{Builder\ coordonnée},
puis regroupé l'ensemble dans une collection Postman de \textbf{872 requêtes PATCH}, organisées
en cinq lots correspondant aux cinq fichiers d'import de couples. Les cinq lots ont été exécutés
après la création des couples, et l'intégralité des requêtes est passée sans erreur.

Ce contournement m'a appris quelque chose que je n'avais pas formulé aussi clairement
auparavant. Dans un projet de reprise, la limite d'un format d'import n'est pas nécessairement
une limite du système lui-même. Le fichier plat et l'API sont deux portes d'entrée vers le même
modèle de données, et savoir basculer de l'une à l'autre au bon moment, plutôt que d'accepter
une perte de données ou de retomber dans la saisie manuelle, est une compétence que la reprise
du premier semestre ne m'avait pas donné l'occasion de mobiliser.

\subsubsection{Traitement des cas particuliers}

Une chaîne pilotée par formules ne traite pas tout. Neuf tiers ont demandé une intervention
manuelle, répartis en quatre familles : trois adresses situées hors de France, que SPO refusait
de créer avec un pays étranger et qui ont été créées en France puis rectifiées ; trois tiers
dont la nature était mal renseignée dans Vadimm, corrigée dans la feuille pivot avant génération
puisque la nature commande la clé d'unicité et la feuille de sortie ; un tiers relevant d'une
opération de dation, déjà présent dans SPO ; et deux doublons portant deux adresses divergentes.

Ce dernier cas est le seul qui ait demandé un arbitrage véritable. La déduplication savait dire
qu'il s'agissait de la même personne, elle ne pouvait pas dire où cette personne habite
aujourd'hui. Après vérification des dates dans Vadimm, nous avons conservé l'adresse rattachée à
la vente la plus récente.

\subsection{Bilan de l'import}

L'intégralité des imports CSV et des requêtes PATCH a été exécutée avec succès.

\begin{longtable}[]{@{}
  >{\raggedright\arraybackslash}p{(\columnwidth - 4\tabcolsep) * \real{0.6364}}
  >{\raggedleft\arraybackslash}p{(\columnwidth - 4\tabcolsep) * \real{0.1212}}
  >{\raggedright\arraybackslash}p{(\columnwidth - 4\tabcolsep) * \real{0.2424}}@{}}
\toprule\noalign{}
\begin{minipage}[b]{\linewidth}\raggedright
Objet créé dans SPO
\end{minipage} & \begin{minipage}[b]{\linewidth}\raggedleft
Volume
\end{minipage} & \begin{minipage}[b]{\linewidth}\raggedright
Canal
\end{minipage} \\
\midrule\noalign{}
\endhead
\bottomrule\noalign{}
\endlastfoot
Tiers & 1~434 & Import CSV \\
Adresses & 553 & Import CSV \\
Courriels (personnes physiques et morales) & 469 & Import CSV \\
Téléphones (personnes physiques et morales) & 538 & Import CSV \\
Représentants (couples et sociétés) & 1~871 & Import CSV \\
Rôles « Client » & 1~434 & Import CSV \\
Coordonnées des couples & 872 & Requêtes PATCH \\
\end{longtable}

Les 1~434 tiers se répartissent en \textbf{881 couples}, \textbf{462 personnes physiques} et
\textbf{91 personnes morales}, répartition que représente la figure~\ref{fig:natures}. Cette
structure justifie après coup l'attention portée aux couples, qui représentent plus de six
tiers sur dix et concentrent la contrainte de format qui a imposé le recours à l'API.

\begin{figure}[!htbp]
\centering
\includegraphics[width=0.65\textwidth]{fig_2_4_repartition_natures.png}
\caption{Répartition des 1~434 tiers créés selon leur nature}
\label{fig:natures}
\end{figure}

\section{Difficultés, limites et perspectives}

Cette mission a comporté sa part de difficultés, dont trois méritent d'être soulignées.

La première, et la plus fréquente, tient au fait qu'un système imposait une règle qui n'était
écrite nulle part. C'est vrai de l'ordre d'injection, du comportement des identifiants ou du
refus des adresses étrangères. Aucune de ces règles ne figurait dans la documentation de
l'éditeur, et je les ai toutes découvertes en les heurtant, par lecture des codes d'erreur. Cela
m'a conduite à documenter systématiquement ce que je découvrais, puisque j'étais la seule à en
avoir la trace.

La deuxième est méthodologique et concerne la définition des clés d'unicité. Ce point n'a rien
de technique : décider que deux lignes désignent la même personne relève d'une règle métier, que
personne ne m'avait formulée et qu'il a fallu construire. J'ai pu m'appuyer sur les échanges avec
les équipes commerciales pour la valider, mais c'est la décision qui engage le plus dans toute
la mission.

La troisième tient à la méthode elle-même, et je préfère la nommer que la passer sous silence.
La manière dont j'ai vérifié mon travail consiste à s'assurer que ce qui est arrivé dans SPO
correspond à ce qui est parti de Vadimm, par paliers de validation puis par comptage à chaque
étage de l'entonnoir. Cette vérification prouve qu'une donnée est bien arrivée, elle ne prouve
pas qu'elle est juste. Si une clé d'unicité était mal choisie, l'erreur passerait le contrôle
sans être vue, puisque les compteurs tomberaient juste. Je reviens sur ce point au chapitre~5.

Au-delà de ces difficultés, cette mission ouvre des perspectives. Le référentiel commercial est
aujourd'hui reconstitué selon des règles écrites et rejouables, mais le chantier n'est pas clos.
La reprise des historiques budgétaires reste devant nous, et c'est elle qui conditionne la suite
du projet du second semestre, comme je l'expose au chapitre~3. Reste également à construire la
synchronisation permanente entre l'ERP et le CRM : la reprise des tiers a créé le référentiel,
elle n'empêche pas encore la double saisie.
